% Part: second-order-logic
% Chapter: metatheory
% Section: compactness

\documentclass[../../../include/open-logic-section]{subfiles}

\begin{document}

\olfileid{sol}{met}{com}

\olsection{द्वितीय-क्रम तर्कशास्त्र संहत नाही}

\begin{explain}
वाक्यांचा संच~$\Gamma$ याचा प्रत्येक सांत उपसंच
पूर्ततायोग्य असेल, तर त्याला \emph{सांततः पूर्ततायोग्य}
म्हणतो. प्रथम-क्रम तर्कशास्त्राचा असा गुणधर्म आहे की
!!{sentence}s चा संच~$\Gamma$ सांततः पूर्ततायोग्य
असेल, तर तो पूर्ततायोग्य असतो. या गुणधर्माला
\emph{संहतता} म्हणतात. याचे परिणाम-संबंध वापरणारे
समतुल्य रूप असे: $\Gamma \Entails !A$ असेल,
तर काही सांत उपसंच~$\Gamma_0 \subseteq \Gamma$
यासाठी आधीच $\Gamma_0 \Entails
!A$ असते. या रूपात हा संपूर्णता प्रमेयाचा थेट
परिणाम आहे: $\Gamma \Entails !A$ असेल, तर
संपूर्णतेमुळे $\Gamma \Proves !A$. पण
!!{derivation} मध्ये~$\Gamma$ मधील सांत
संख्येनेच !!{sentence}s वापरता येतात.

द्वितीय-क्रम तर्कशास्त्रासाठी संहतता लागू होत नाही.
द्वितीय-क्रम !!{sentence}s चे असे संच आहेत की ते
सांततः पूर्ततायोग्य असतात, पण पूर्ततायोग्य नसतात;
आणि ते एखाद्या~$!A$ चे परिणाम देतात, पण त्यांचा
कोणताही सांत उपसंच त्या~$!A$ चे परिणाम देत नाही.
\end{explain}


\begin{thm}
\ollabel{thm:sol-not-compact}
द्वितीय-क्रम तर्कशास्त्र संहत नाही.
\end{thm}

\begin{proof}
हे आठवा:
\[
\fn{Inf} \ident \lexists[u][(\lforall[x][\lforall[y][(\eq[u(x)][u(y)]
      \lif \eq[x][y])]] \land \lexists[y][\lforall[x][\eq/[y][u(x)]]])]
\]
!!a{structure} चा प्रांत अनंत असेल, जर आणि फक्त जर
तिच्यात या सूत्राची पूर्ति होते. प्रांतात किमान
$n$ !!{element}s आहेत, असे सांगणारे वाक्य
$!A^{\ge n}$ असो; उदाहरणार्थ,
\[
!A^{\ge n} \ident \lexists[x_1][\dots\lexists[x_n][(\eq/[x_1][x_2]
    \land \eq/[x_1][x_3] \land \dots \land \eq/[x_{n-1}][x_n])]].
\]
!!{sentence}s चा पुढील संच पाहू:
\[
\Gamma = \{\lnot \fn{Inf}, !A^{\ge 1}, !A^{\ge 2}, !A^{\ge 3}, \dots\}.
\]
तो सांततः पूर्ततायोग्य आहे: कोणताही सांत उपसंच~$\Gamma_0
\subseteq \Gamma$ घेतल्यास असा $k$ निवडता येतो की
$!A^{\ge k} \in \Gamma$ असते, पण $n>k$ साठी
कोणतेही $!A^{\ge n} \in \Gamma_0$ नसते.
$\Domain{M}$ मध्ये $k$ !!{element}s असतील, तर
$\Sat{M}{\Gamma_0}$. मात्र $\Gamma$ पूर्ततायोग्य
नाही: $\Sat{M}{\lnot \fn{Inf}}$ असेल, तर
$\Domain{M}$ सांत असला पाहिजे; समजा त्याचा
आकार~$k$ आहे. मग $\Sat/{M}{!A^{\ge k+1}}$.
\end{proof}

\begin{prob}
संच~$\Gamma$ आणि !!a{sentence}~$!A$ यांचे असे
उदाहरण द्या की $\Gamma \Entails !A$, पण प्रत्येक
सांत उपसंच~$\Gamma_0 \subseteq
\Gamma$ साठी $\Gamma_0 \Entails/ !A$.
\end{prob}


\end{document}
